\section{Field arithmetic: two representations, one model}
\label{sec:arith}

A step is dominated by field multiplication: two products for the
addition and, with $W$ walks sharing one inversion by Montgomery's
trick~\cite{montgomery1987}, $3(W-1)/W$ more for the product tree,
plus $1/W$ of an Itoh--Tsujii inversion~\cite{itoh1988} that costs
eight multiplications in a basis where squaring is free (the client's
addition chain is $1,2,4,8,16,32,64,65,130$, the FPGA's
$1,2,4,\ldots,128,130$).\art{ecc2k130/include/packed131.h}\art{hdl/ecc2k130/README.md}
Going from $W=4$ to $W=32$ is worth $60\%$ on the CPU client because
it drives the amortised inversion from $25\%$ to $3\%$ of the
multiplication budget.\art{ecc2k130/README.md, ``Throughput''} The
representation of a multiplication is therefore the design decision,
and the repository carries two.

\subsection{Bitsliced: 32 walks in the bits of a word}
\label{sec:arith-bitsliced}

The original client is bitsliced in the sense of
Biham~\cite{biham1997}: one word operation advances $W$ walks at once,
one per bit lane, and a field operation is a straight-line program of
word-wide Boolean instructions generated from the field model. The
device uses 32-bit words; the host picks the widest it has (AVX-512
gives 512 lanes, AVX2 256, NEON 128), and the arithmetic source is
identical at every width.\art{ecc2k130/README.md, ``Word width''}

Multiplication follows Bernstein and Lange~\cite{bernsteinlange2010onb}:
convert both operands from the normal basis to the optimal polynomial
basis $\{1,c,c^2,\ldots\}$ with $c=\gamma_1$, multiply as polynomials,
and convert the 261-coefficient product back. Both conversions are
generated from the substitution $c=z+1/z$, which splits by parity as
$H(c)=H_{\mathrm{even}}(c^2)+c\,H_{\mathrm{odd}}(c^2)$; since squaring
is an index permutation this gives an $O(K\log K)$ conversion. The
generated \texttt{multPrep} is 331 instructions and the double-size
inverse conversion 903, against 325 and 909 for the hand-designed
chains in the literature.\art{ecc2k130/README.md, ``Findings''}

Two findings from this backend change the standard accounting. First,
three-input logic fusion (\texttt{LOP3} on NVIDIA, \texttt{vpternlogd}
on AVX-512, \texttt{BSL}/\texttt{EOR3} on Arm with \texttt{FEAT\_SHA3})
is worth $1.65\times$: a 131-bit multiplication is $14{,}554$ two-input
bit operations but $8{,}859$ instructions once $x\oplus y\oplus z$ and
$x\oplus(y\wedge z)$ are single instructions. Because the fused form
charges nothing for the XOR that accompanies an AND, schoolbook leaves
beat Karatsuba leaves~\cite{karatsuba1962}: full recursion costs
$18{,}819$ bit operations and $14{,}770$ instructions, while stopping at
a 17-bit schoolbook leaf costs more bit operations and fewer
instructions. A search over balanced and unbalanced splits found nothing below
$8{,}859$; Toom-3 over $\FF_2$, the source of Bernstein's
$11{,}961$-bit-operation chain~\cite{bernstein2009batch}, is not
implemented.\art{ecc2k130/README.md, ``What is not done''} Second, code size rather than arithmetic is the wall:
inlining everything produced $285{,}096$ PTX instructions for one
kernel and \texttt{ptxas} had not finished register allocation after
five minutes; marking the multiplication, inversion, conversions,
leaf, weight tree and $\sigma^j$ as real subroutines brings it to
$36{,}553$ lines and a $2.4$~s assembly.\art{ecc2k130/README.md, ``Findings''}

The Hamming-weight tree is 268 instructions against 654 bit
operations, because a full adder is two \texttt{LOP3}s. The whole bitsliced iteration is about $60{,}000$ instructions per 32
walks on $\mathrm{GF}(2^{131})$, a static count; on the GPU that is
the $0.85$~B/s control of Section~\ref{sec:ledger}, and the
instruction-count ceiling of the later packed kernel on the
RTX~PRO~6000 is $17.6$~B/s ($27.1$~T lane-operations/s over about
$1{,}540$ lane-operations per update, a different
count).\art{ecc2k130/THROUGHPUT-CEILING.md}

\subsection{Packed: one walk per thread, carry-less products}
\label{sec:arith-packed}

The packed backend holds one walk's coordinates in five 32-bit words
per thread and multiplies with carry-less
instructions.\art{ecc2k130/PACKED.md} Coordinates live in
$\FF_2[c]/(m)$ with $c=\gamma_1$, so a 131-bit product is a Karatsuba
over carry-less multiplies, six \texttt{CLMAD}, followed by a
reduction;\art{ecc2k130/ITERATION-FUNCTION.md, §1} the normal basis,
where the weight, the branch selection and the distinguished-point
test live, is reached through the same Bernstein--Lange conversions,
realised as shift-and-mask networks and, for the Frobenius powers, as
Bene\v{s} permutation networks.\art{ecc2k130/FROBENIUS-NETWORK.md}
Before CUDA~13.3 the carry-less product was built from integer
multiplies; the native \texttt{clmad.lo}/\texttt{hi} instructions on
Blackwell gave $+22.4\%$ in Section~\ref{sec:ledger} and are the
reason the campaign preset requires that
compiler.\art{ecc2k130/NATIVE-CARRYLESS.md}

The accounting unit changes with the representation. In the packed kernel the shipping walk issues $4.81$ products per
update at batch 16 and about $2{,}324$ static ALU
slots;\art{ecc2k130/ITERATION-FUNCTION.md, §1 and §3} the step count
that matters is the \texttt{CLMAD} count, and Section~\ref{sec:floor}
prices it. The reduction \texttt{reducePolynomial131} alone is $440.6$
lane-instructions per update, $29\%$ of the integer work of the
$20$~B/s build.\art{ecc2k130/ROOFLINE.md}

\subsection{Both at once, and bit-for-bit}
\label{sec:arith-compat}

The two backends, the CPU client at every word width, the VHDL
engine and the sibling library's imported kernel are kept bit-for-bit
compatible: a distinguished-point record written by any of them
re-walks under any other, and every GPU report is re-walked on the
host reference before a rate is recorded. The host check re-slices
packed words into the same permuted type-II basis and compares
vectors, not classes.\art{ecc2k130/README.md, ``Validation''} The exceptional case $\sigma^j(R)=\pm R$, which makes the addition's
denominator zero, has probability about $2^{-131}$ per step. The
$\sigma$-walk client and the VHDL engine do not special-case it: a
zero denominator inverts to zero and the batch's points become
garbage, which the next distinguished-point check or the guard
discards. The sibling C model and the table walk abort the step
instead.\art{ecc2k130/include/ref.h}\art{hdl/ecc2k130/README.md}
